\chapter{Scénáře uživatelského testování}


\section{Zadání testování}

\subsection{Scénář 1: Vyhledání obchodu pomocí filtrování}
\begin{description}
    \item[Zadání:] Najděte obchod v okolí do 5 km, který prodává zeleninu a má hodnocení alespoň 3 hvězdičky.
    \item[Upřesnění:] Využijte filtry pro kategorii, vzdálenost a hodnocení.
\end{description}

\subsection{Scénář 2: Vytvoření obchodu}
\begin{description}
    \item[Zadání:] Zaregistrujte se a vytvořte obchod zaměřený na prodej domácích vajec a brambor.
    \item[Upřesnění:] Přidejte alespoň 2 fotografie (1 fotografie místa prodeje, 1 fotografie produktu), u produktů vyplňte alespoň název a cenu za určité množství a nastavte otevírací dobu (např.\ formou textové informace pro kontaktování).
\end{description}

\subsection{Scénář 3: Úprava dostupnosti produktů}
\begin{description}
    \item[Zadání:] Ve vytvořeném obchodě nastavte, že produkt \enquote{brambory} není aktuálně dostupný.
    \item[Upřesnění:] Změňte stav produktu na \enquote{není skladem}.
\end{description}

\subsection{Scénář 4: Úprava uživatelského profilu}
\begin{description}
    \item[Zadání:] Upravte svůj uživatelský profil.
    \item[Upřesnění:] Přidejte profilovou fotografii, vyplňte jméno a přidejte alespoň jeden preferovaný kontakt.
\end{description}

\subsection{Scénář 5: Vytvoření recenze}
\begin{description}
    \item[Zadání:] Vytvořte recenzi u libovolného obchodu.
    \item[Upřesnění:] Přidejte hodnocení i text recenze.
\end{description}